\documentclass{beamer}					% Document class

\usepackage[english]{babel}				% Set language
\usepackage[utf8x]{inputenc}			% Set encoding

% 加入 listings 套件
\usepackage{listings}
\usepackage{xcolor}

% 設定 listings 樣式
\lstset{
    language=R,
    basicstyle=\ttfamily\small,
    backgroundcolor=\color{gray!10},
    frame=single,
    breaklines=true,
    showstringspaces=false,
    keywordstyle=\color{blue},
    commentstyle=\color{green!50!black},
    stringstyle=\color{red},
    numbers=none,
    xleftmargin=10pt,
    xrightmargin=10pt,
    framexleftmargin=5pt,
    framexrightmargin=5pt
}

\mode<presentation>						% Set options
{
  \usetheme{default}					% Set theme
  \usecolortheme{default} 				% Set colors
  \usefonttheme{default}  				% Set font theme
  \setbeamertemplate{caption}[numbered]	% Set caption to be numbered
}

\usepackage{graphicx}					% For including figures
\usepackage{booktabs}					% For table rules
\usepackage{hyperref}					% For cross-referencing
\usepackage{subcaption}  % 記得加在 preamble
\usepackage{multicol}

\title{Data Type \& Data Structure}	% Presentation title
\author{Rex Shao-Yu Wang}								% Presentation author
\institute{Department of Agronomy, National Taiwan University}					% Author affiliation
\date{\today}									% Today's date	

\begin{document}

% Title page
\begin{frame}
  \titlepage
\end{frame}

% Outline
\begin{frame}{Outline}
  \tableofcontents
\end{frame}

\section{Data Type}
\begin{frame}{Basic Data Type}
\begin{tabular}{lll}
\toprule
\textbf{Type} & \textbf{Range} & \textbf{Assign a Value} \\
\midrule
Numeric  & Numbers & x \textless- 1 \\
Character  & Text strings  & my\_name \textless- "Rex" \\
Logical & TRUE or FALSE  & Stevenistheteacher \textless- TRUE \\
\bottomrule
\end{tabular} \\
\vspace{20pt}
\texttt{class()}: To check the data type
\end{frame}

\section{Data Structure}
\begin{frame}{Basic Data Structure}
\begin{tabular}{lll}
\toprule
\textbf{Type} & \textbf{Range} \\
\midrule
Vector   & A list of items that are of the same type \\
Matrix   & A two-dimensional data set with columns and rows \\
Data.Frame  & Data displayed in a format as a table \\
\bottomrule
\end{tabular} 

\begin{figure}
    \centering
    \includegraphics[width=0.7\linewidth]{Data_types_structure_image.jpg}
\end{figure}
\end{frame}

\begin{frame}[fragile]{Vectors}  % 注意：加入 fragile 選項很重要！
\texttt{c()} stands for concatenate, which combines elements into a vector
\begin{lstlisting}
score <- c(60, 75, 89, 95, 71, 90)
name <- c("John", "James", "Hank", "Jeff", "Ryan", "Ray")
class <- c(1, 1, 1, 2, 2, 2)
\end{lstlisting}

\texttt{rep()}: Replicate elements of vectors and lists
\begin{lstlisting}
rep(x = 1, times = 9)
rep(x = c(1, 2), each = 3) 
rep(x = c(1, 2, 3), times = c(2, 3, 5))
\end{lstlisting}

\texttt{seq()}: Tool used to generate regular sequences of numbers
\begin{lstlisting}
seq(from = 1, to = 5, by = 1)
seq(1, 9, by = 2)   
seq(0, 1, length.out = 11)
\end{lstlisting}
\end{frame}

\begin{frame}[fragile]{Vectors}
What will these commands do?
\begin{lstlisting}
score[1]
score[1:3]
score[-1]
\end{lstlisting}
\begin{figure}
    \centering
    \includegraphics[width=0.7\linewidth]{截圖 2026-03-07 凌晨2.20.14.png}
\end{figure}

typeof() determines the type or storage mode of any object
\begin{lstlisting}
typeof(score)
typeof(name)
\end{lstlisting}

You can also use length() to know the number of elements in the vector
\begin{lstlisting}
length(score)
length(name)
\end{lstlisting}
\end{frame}

\begin{frame}[fragile]{Exercises: Vectors in R}
\begin{enumerate}
\item Use the seq() function to generate the sequence \\  9, 18, 27, 36, 45.
\item Use the rep() function to generate the sequence \\ 1 1 1 1 1 2 2 2 2 3 3 3 4 4 5
\item Given the following vectors:
\begin{lstlisting}
students <- c("Alice", "Bob", "Chris", "David", "Eve")
\end{lstlisting}
Write the R commands to:
\begin{enumerate}
    \item Extract the first element from the \texttt{students} vector.
    \item Extract the first three elements from the \texttt{students} vector.
    \item Extract all names from \texttt{students} \textbf{except} for the second one ("Bob").
\end{enumerate}

\end{enumerate}
\end{frame}

\begin{frame}[fragile]{Matrices}
\begin{columns}

\column{0.7\textwidth}
\texttt{matrix(data, nrow, ncol, byrow)} creates a matrix
\begin{itemize}
    \item data: take a set of values
    \item nrow: specifies the number of rows
    \item ncol: specifies the number of columns
    \item byrow: indicates that the matrix is filled row-wise

\end{itemize}
\begin{lstlisting}
height <- c(175, 168, 180, 154, 185, 162, 170, 187, 149)
matrix(data = height, nrow = 3, ncol = 3, byrow = T)
matrix(data = height, nrow = 3, ncol = 3, byrow = F)
\end{lstlisting}

\column{0.3\textwidth}
\vspace{0}
\begin{figure}
    \centering
    \includegraphics[width=1.15\linewidth]{截圖 2026-03-07 上午10.14.08.png}
\end{figure}

\begin{figure}
    \centering
    \includegraphics[width=1.15\linewidth]{Representation-of-confusion-matrix-True-Positive-True-positive-shows-the-number-of.png}
\end{figure}


\end{columns}
\end{frame}

\begin{frame}[fragile]{Matrices}

\begin{columns}
\column{0.55\textwidth}
What will these commands do?
\begin{lstlisting}
height.mx[1]
height.mx[1, 2]
height.mx[2, -1]
\end{lstlisting}
Check the structure with
\begin{lstlisting}
dim(height.mx)
length(height.mx)
nrow(height.mx)
ncol(height.mx)
\end{lstlisting}

\column{0.45\textwidth}
\begin{figure}
    \centering
    \includegraphics[width=1\linewidth]{截圖 2026-03-07 上午10.21.59.png}
\end{figure}
\end{columns}
\end{frame}

\begin{frame}[fragile]{Data Frames}
\texttt{data.frame()} creates data frames, coupled collections of variables
\begin{columns}
\column{0.55\textwidth}
\begin{lstlisting}
information <- data.frame(
  Name = name,
  Class = class,
  Score = score
)
\end{lstlisting}

\column{0.45\textwidth}
\begin{figure}
    \centering
    \includegraphics[width=1.1\linewidth]{data-frames-in-r-thumbnail.png}
\end{figure}
\end{columns}
\end{frame}

\begin{frame}[fragile]{Data Frame}
\begin{columns}
\column{0.55\textwidth}
What will these commands do?
\begin{lstlisting}
information["Name"]
information[2]
information[2, ]
information[1:3]
information$Name
\end{lstlisting}
Check the structure with
\begin{lstlisting}
colnames(information)
rownames(information)
length(information)
nrow(information)
ncol(information)
\end{lstlisting}

\column{0.45\textwidth}
\begin{figure}
    \centering
    \includegraphics[width=1.1\linewidth]{截圖 2026-03-09 上午10.16.03.png}
\end{figure}
\end{columns}
\end{frame}

\section{Lab Report}
\begin{frame}{Quarto}
A framework enables you to weave together content and executable code into a finished document\\
\vspace{6}
Mostly contains three parts
\begin{itemize}
\item YAML header
\item Chunks of R code
\item Text mixed with simple text
\end{itemize}
\begin{figure}
    \centering
    \includegraphics[width=0.9\linewidth]{截圖 2026-03-07 凌晨1.32.19.png}
\end{figure}    
\end{frame}

\begin{frame}{Lab Report}
Please make sure your R and RStudio are the latest version\\
Otherwise, Quarto might not work
\begin{itemize}
    \item The report file will be in the NTU COOL Assignment section
    \item Submit your file in pdf format
    \item Deadline: 3/15 (Sun) 23:59
    \item Late submission is not accepted
\end{itemize}
    
\end{frame}

\begin{frame}{Lab Report}
Download the file in NTU COOL
\begin{figure}
    \centering
    \includegraphics[width=0.9\linewidth]{截圖 2026-03-07 凌晨1.16.30.png}
\end{figure}
Type your student id and name in the author section in YAML    
\begin{figure}
    \centering
    
\includegraphics[width=0.9\linewidth]{截圖 2026-03-07 凌晨1.46.41.png}
\end{figure}
\end{frame}

\begin{frame}{Lab Report}
Type in your command in the chunk. You can check your answer via the green arrow
\begin{figure}
    \centering
    \includegraphics[width=0.9\linewidth]{截圖 2026-03-07 凌晨1.49.35.png}
\end{figure}

After you've finished all the tasks, click render on the tool list
\begin{figure}
    \centering
    \includegraphics[width=0.9\linewidth]{截圖 2026-03-07 凌晨1.50.33.png}
\end{figure}
\end{frame}

\begin{frame}{Lab Report}
Download the pdf file and submit it to NTU COOL\\
\begin{figure}
    \centering
    \includegraphics[width=1\linewidth]{截圖 2026-03-07 凌晨1.53.38.png}
\end{figure}
*For those who is unable to do it on Quarto, you can also take a screenshot of your code and output, and paste it on a pdf file
\end{frame}
\end{document}
